\section{总结与延伸}

本期的核心收获：

\begin{enumerate}
  \item \textbf{等效宽度} $k \times m = 6144$ 是理解 MoE 计算量的关键——它等效于一个 $3 \times d_{\text{model}}$ 宽的 Dense FFN
  \item 加权求和与 concat 在 FLOPs 和矩阵 rank 上等价，MoE 通过多个「瘦」Expert 承载多样性
  \item Attention 是 token 间的因果交互，FFN/MoE 是 token 内的独立映射——两者角色本质不同
  \item 无论 Dense 还是 MoE，FFN 部分始终是参数量的绝对大户
  \item RMSNorm 每层 2 个 + 末尾 1 个，参数量虽小但概念上不可忽略
\end{enumerate}

\subsection{拓展阅读}

\begin{itemize}
  \item HuggingFace Transformers 源码：\texttt{models/qwen3\_moe/}
  \item Switch Transformer (Fedus et al., 2021)：MoE 在 Transformer 中的经典应用
  \item DeepSeek MoE 技术报告：更细粒度的 Expert 设计
\end{itemize}

\end{document}
